% ch5_c (书页 217-228 / p232-p243)
%--B-TAIL--
%JOIN%解，因而应当预言一种在时空小区域内退化为等效原理（EP）预言的红移。让我们来看看这是如何实现的。

% ---- 书页 217 / p232 ----
考虑一个四速度为 $U^\mu$、在 Schwarzschild 坐标中静止（$U^i = 0$）的观者。我们本可以允许这个观者运动，但那只不过是在引力效应之上叠加一个通常的多普勒频移而已。四速度满足 $U_\mu U^\mu = -1$，对于 Schwarzschild 时空中的静止观者，这就意味着
\begin{equation*}
U^0 = \left(1 - \frac{2GM}{r}\right)^{-1/2}.
\tag{5.99}
\end{equation*}
任何一个这样的观者，测量沿类光测地线 $x^\mu(\lambda)$ 传播的光子的频率，得到的都是
\begin{equation*}
\omega = -g_{\mu\nu}U^\mu\frac{\mathrm{d}x^\nu}{\mathrm{d}\lambda}.
\tag{5.100}
\end{equation*}
实际上，正是这个关系定义了 $\lambda$ 的归一化。于是我们有
\begin{align*}
\omega &= \left(1 - \frac{2GM}{r}\right)^{1/2}\frac{\mathrm{d}t}{\mathrm{d}\lambda}
\tag{5.101}\\
&= \left(1 - \frac{2GM}{r}\right)^{-1/2}E,
\tag{5.102}
\end{align*}
其中 $E$ 由式 (5.61) 定义，只是应用于光子的轨迹。$E$ 是守恒量，因此在不同的径向距离处测量时，$\omega$ 显然会取不同的值。对于在 $r_1$ 处发射、在 $r_2$ 处被观测到的光子，观测到的频率之间将有如下关系
\begin{equation*}
\frac{\omega_2}{\omega_1} = \left(\frac{1 - 2GM/r_1}{1 - 2GM/r_2}\right)^{1/2}.
\tag{5.103}
\end{equation*}
这是频移的一个精确结果；在 $r \gg 2GM$ 的极限下，我们得到
\begin{align*}
\frac{\omega_2}{\omega_1} &= 1 - \frac{GM}{r_1} + \frac{GM}{r_2}\\
&= 1 + \Phi_1 - \Phi_2,
\tag{5.104}
\end{align*}
其中 $\Phi = -GM/r$ 是 Newton 引力势。这告诉我们，当 $\Phi$ 增大时频率会降低——而这正是我们从引力场中向外爬升时发生的事情：于是就出现了红移。（落向引力体的光子则会发生蓝移。）我们看到，$r \gg 2GM$ 的结果与基于等效原理的计算相一致。

引力红移最早由 Pound 与 Rebka 于 1960 年探测到，他们让 $\gamma$ 射线向上穿行仅仅 72 英尺（哈佛大学物理系大楼的高度）的距离。后来的检验变得越来越精确，往往利用人造飞船或放置在飞机上的原子钟。在所有这些情形中，与 Einstein 预言的符合程度都极为出色。

% ---- 书页 218 / p233 ----
自 Einstein 提出三大经典检验以来，人们又提出了对广义相对论的若干进一步检验。其中最著名的当然是脉冲双星，我们将在第 7 章讨论。另一个是引力时间延迟，由 Shapiro 发现并观测到，同样在第 7 章讨论。在一个非常不同的情境下，大爆炸核合成为广义相对论提供了一个宇宙学检验——它发生在宇宙仅诞生几秒钟的年代，将在第 8 章讨论。现代的进展还引入了大量新的检验；全面的介绍可参阅 Will (1981)。

\section{Schwarzschild 黑洞}

现在，对于麻烦半径 $r = 2GM$ 之外测地线的行为我们已经有所了解——这正是太阳系以及其他大多数天体物理情形所关心的区域。接下来我们转而研究这样一些天体：即使在半径小于 $2GM$ 处，它们仍然由 Schwarzschild 解描述——黑洞。（我们暂时先使用“黑洞”这个术语，尽管还没有为这样一种天体引入确切的含义。）

理解一个时空的几何的办法之一，是考察由光锥定义的因果结构。因此我们考虑径向类光曲线，即 $\theta$ 和 $\phi$ 保持不变且 $\mathrm{d}s^2 = 0$ 的曲线：
\begin{equation*}
\mathrm{d}s^2 = 0 = -\left(1 - \frac{2GM}{r}\right)\mathrm{d}t^2 + \left(1 - \frac{2GM}{r}\right)^{-1}\mathrm{d}r^2,
\tag{5.105}
\end{equation*}
由此我们得到
\begin{equation*}
\frac{\mathrm{d}t}{\mathrm{d}r} = \pm\left(1 - \frac{2GM}{r}\right)^{-1}.
\tag{5.106}
\end{equation*}
当然，这量度的正是 $t$-$r$ 平面的时空图上光锥的斜率。$r$ 很大时斜率为 $\pm 1$，与平直空间中一样；而当我们接近 $r = 2GM$ 时，$\mathrm{d}t/\mathrm{d}r \to \pm\infty$，光锥便“闭合”起来，如图 5.7 所示。因此，接近 $r = 2GM$ 的光线似乎永远到不了那里——至少在这个坐标系中是如此；它看起来是渐近地趋于这个半径。

%FIG{5.7}: {在 Schwarzschild 坐标中，当我们接近 $r = 2GM$ 时，光锥看似逐渐闭合。}
正如我们将要看到的，到不了 $r = 2GM$ 只是一种表象，光线（或者有质量的粒子）实际上到达这个半径毫无困难。但远处的观者永远无法知晓这一点。如果我们待在外面，而一位勇敢的从事观测的广义相对论学家跳入黑洞、并不停地发回信号，我们只会看到信号到达我们这里越来越慢，如图 5.8 所描绘。在习题中会要求你更仔细地考察这一现象。当落入的观者接近 $r = 2GM$ 时，其固有时中任何固定的间隔 $\Delta\tau_1$，从我们的观点来看都对应于一个越来越长的间隔 $\Delta\tau_2$。这种情形永远持续下去；我们永远不会看到
% ---- 书页 219 / p234 ----
观者穿过 $r = 2GM$，只会看到他们运动得越来越慢（而且变得越来越红，仿佛因为干了跳进黑洞这么愚蠢的事而感到无地自容）。

我们永远看不到落入的观者到达 $r = 2GM$，这一事实是一个有意义的陈述；但他们在 $t$-$r$ 平面中的轨迹永远到不了那里，这一事实则不是。后者高度依赖于我们的坐标系，而我们想问的是一个更不依赖坐标的问题（比如：“观者能否在自己的固有时有限量之内到达这个半径？”）。做到这一点的最好办法，是换成一组在 $r = 2GM$ 处表现更好的坐标系。现在我们就来寻找一组合适的这样的坐标。当然，坐标变换是无法“推导”出来的，我们只能径直说出新坐标是什么，然后代入公式。不过我们将分几步来建立这些坐标，希望这些选择看起来多少有些缘由。

%FIG{5.8}: {自由落入黑洞的一个信标以恒定固有时间隔 $\Delta\tau_1$ 发出信号。固定 $r$ 处的观者收到信号的时间间隔 $\Delta\tau_2$ 依次变长。}

% ---- 书页 220 / p235 ----
我们现有坐标的问题在于，沿接近 $r = 2GM$ 的径向类光测地线，$\mathrm{d}t/\mathrm{d}r \to \infty$；$r$ 方向上的推进相对于坐标时 $t$ 变得越来越慢。我们可以试着解决这个问题，办法是用一个沿类光测地线变化得更慢的坐标来取代 $t$。首先注意到，我们可以显式地解出刻画径向类光曲线的条件 (5.106)，得到
\begin{equation*}
t = \pm r^* + \text{constant},
\tag{5.107}
\end{equation*}
其中\textbf{乌龟坐标}（tortoise coordinate）$r^*$ 定义为
\begin{equation*}
r^* = r + 2GM\ln\left(\frac{r}{2GM} - 1\right).
\tag{5.108}
\end{equation*}
（乌龟坐标只有在 $r \geq 2GM$ 时才与 $r$ 有合理的关系，但反正在那个范围之外我们的坐标本来也不怎么好。）用乌龟坐标表示，Schwarzschild 度规变为
\begin{equation*}
\mathrm{d}s^2 = \left(1 - \frac{2GM}{r}\right)\left(-\mathrm{d}t^2 + \mathrm{d}r^{*2}\right) + r^2\,\mathrm{d}\Omega^2,
\tag{5.109}
\end{equation*}
其中 $r$ 被看作 $r^*$ 的函数。这代表了一些进展，因为如图 5.9 所示，光锥现在看起来不再闭合；此外，在 $r = 2GM$ 处没有任何度规系数变为无穷（尽管 $g_{tt}$ 和 $g_{r^* r^*}$ 都变为零）。然而我们付出的代价是，感兴趣的曲面 $r = 2GM$ 只是被推到了无限远。

我们的下一步行动，是定义自然地贴合类光测地线的坐标。若令
\begin{align*}
v &= t + r^*\\
u &= t - r^*,
\tag{5.110}
\end{align*}

%FIG{5.9}: {乌龟坐标（式 (5.109)）下的 Schwarzschild 光锥。光锥保持非退化，但曲面 $r = 2GM$ 已被推到无限远。}

% ---- 书页 221 / p236 ----
则落入的径向类光测地线由 $v =$ 常数刻画，而出射的径向类光测地线满足 $u =$ 常数。现在考虑回到原来的径向坐标 $r$，但用新坐标 $v$ 取代类时坐标 $t$。这组坐标就是所谓的\textbf{Eddington--Finkelstein 坐标}。用这组坐标表示，度规为
\begin{equation*}
\mathrm{d}s^2 = -\left(1 - \frac{2GM}{r}\right)\mathrm{d}v^2 + \left(\mathrm{d}v\,\mathrm{d}r + \mathrm{d}r\,\mathrm{d}v\right) + r^2\,\mathrm{d}\Omega^2.
\tag{5.111}
\end{equation*}
这里我们第一次看到了真正进展的迹象。尽管度规系数 $g_{vv}$ 在 $r = 2GM$ 处为零，但并不存在真正的退化；度规的行列式为
\begin{equation*}
g = -r^4\sin^2\theta,
\tag{5.112}
\end{equation*}
它在 $r = 2GM$ 处完全正则。因此度规是可逆的，而且我们一劳永逸地看到：在我们最初的 $(t, r, \theta, \phi)$ 坐标系中，$r = 2GM$ 只不过是一个坐标奇点。在 Eddington--Finkelstein 坐标中，径向类光曲线的条件解得
\begin{equation*}
\frac{\mathrm{d}v}{\mathrm{d}r} =
\begin{cases}
0, & \text{（落入）}\\[6pt]
2\left(1 - \dfrac{2GM}{r}\right)^{-1}. & \text{（出射）}
\end{cases}
\tag{5.113}
\end{equation*}
于是我们能够看清发生了什么：在这个坐标系中，光锥在 $r = 2GM$ 处依然表现良好，而且这个曲面位于有限的坐标值处。追踪类光或类时粒子穿过该曲面的路径没有任何问题。另一方面，确实有某种有趣的事情在发生。虽然光锥并不闭合，但它们确实倾倒了下来，使得当 $r < 2GM$ 时，所有指向未来的路径都指向 $r$ 减小的方向，如图 5.10 所示。

%FIG{5.10}: {式 (5.111) 的 $(v, r)$ 坐标下的 Schwarzschild 光锥。在这些坐标中，指向未来的类时路径可以越过 $r = 2GM$。}

% ---- 书页 222 / p237 ----
曲面 $r = 2GM$ 虽然局部完全正则，但在整体上起着有去无回之点的作用——检验粒子一旦没入其内，就永远无法返回。我们把\textbf{事件视界}（event horizon）定义为粒子永远无法穿过它逃逸到无限远的曲面；在 Schwarzschild 时空中，事件视界位于 $r = 2GM$ 处。（这是一个粗略的定义；在下一章中我们会讲得稍微精确一些。）尽管位于固定的径向坐标处，事件视界却是一个类光曲面而非类时曲面，所以真正使得沿向外方向穿过视界成为不可能的，正是时空自身的因果结构。既然任何东西都无法逃出事件视界，我们也就不可能看到它的内部——这便是\textbf{黑洞}（black hole）这一名称的由来。黑洞不过是时空中的一个被事件视界与无限远分隔开的区域。事件视界的概念是一个整体性的概念；视界的位置是关于时空整体的陈述，不是仅知道该处的几何就能确定的。在更一般的时空中，这一点仍将成立。

我们应该提一下黑洞的几个在大众想象中常常被混淆的特点。首先，黑洞的外部几何与我们本会在恒星或行星之外得到的 Schwarzschild 解是同一个解。特别地，黑洞并不比太阳更会把周围的一切吸进去；在远离 $r = 2GM$ 之外的地方，无论引力源是不是黑洞，粒子的行为都完全一样。其次，关于黑洞存在一种误导性的 Newton 式类比。质量为 $M$ 的引力体外距离 $r$ 处的一个粒子，其 Newton 逃逸速度为
\begin{equation*}
v_{\mathrm{esc}} = \sqrt{\frac{2GM}{r}}.
\tag{5.114}
\end{equation*}
如果我们天真地问 Newton 逃逸速度在哪里等于光速，会发现恰好就是 $r = 2GM$。尽管光速在 Newton 理论中并不扮演任何基本角色，下面这件事看上去仍颇具挑衅意味：光——被当作以速度 $c$ 运动的惯性粒子——似乎无法从质量为 $M$、半径小于 $2GM$ 的天体逃出。但这种情况与我们在广义相对论中看到的存在深刻的差别。逃逸速度是粒子想要沿自由轨迹从引力源逃出时最初需要具有的速度；但没有什么能阻止我们考虑加速轨迹——比如，可以想象选取一个加速度，使得粒子以某个恒定的速度稳定地远离大质量天体。因此，假想的 Newton 黑洞将不会具有\emph{任何东西}都无法逃出这一关键性质；而在广义相对论中，任意的类时路径都必须待在各自的光锥之内，因而永远无法逃出事件视界。

\section{最大延拓的 Schwarzschild 解}

让我们回顾一下已经做了什么。出于对我们的坐标也许并不适用于整个流形的怀疑，我们已把原来的坐标 $t$ 换成了新坐标 $v$；后者的一个良好性质是，如果我们沿 $v =$ 常数的径向类光曲线减小 $r$，就会径直穿过事件视界而没有任何问题。

% ---- 书页 223 / p238 ----
的确，真正完成这趟旅行的局部观者未必知道事件视界是在何时被越过的——局部几何与其他任何地方并无不同。因此我们得出结论：我们的怀疑是对的，最初的坐标系未能很好地覆盖整个流形。区域 $r \leq 2GM$ 理所当然应当包含在我们的时空之中，因为物理粒子可以轻松到达并穿过那里。然而，这并不能保证我们已经大功告成；也许我们还能沿其他方向延拓流形。

实际上确实还有其他方向。在 $(v, r)$ 坐标系中，我们能够沿指向未来的路径穿过事件视界，却不能沿指向过去的路径。这显得不合情理，因为我们一开始面对的是一个不依赖时间的解。但我们本可以选择 $u$ 而不是 $v$，那样的话度规将是
\begin{equation*}
\mathrm{d}s^2 = -\left(1 - \frac{2GM}{r}\right)\mathrm{d}u^2 - \left(\mathrm{d}u\,\mathrm{d}r + \mathrm{d}r\,\mathrm{d}u\right) + r^2\,\mathrm{d}\Omega^2.
\tag{5.115}
\end{equation*}
现在我们又能穿过事件视界了，但这一次只能沿指向过去的曲线，如图 5.11 所示。

这也许令人惊讶：我们可以自洽地沿指向未来或指向过去的曲线穿过 $r = 2GM$，但会到达不同的地方。这其实本来就该预料到，因为从定义 (5.110) 看，如果保持 $v$ 不变而减小 $r$，必有 $t \to +\infty$；而如果保持 $u$ 不变而减小 $r$，则有 $t \to -\infty$。（当 $r \to 2GM$ 时，乌龟坐标 $r^* \to -\infty$。）于是我们沿两个不同的方向延拓了时空：一个朝向未来，一个朝向过去。

下一步本来可以沿着类空测地线走，看看是否还会发现更多的区域。答案是肯定的——我们会到达时空的另一块区域；不过让我们抄个近道，直接定义一组处处都好用的坐标。第一个猜测也许是同时使用 $u$ 和 $v$（来代替 $t$ 和 $r$），

%FIG{5.11}: {式 (5.115) 的 $(u, r)$ 坐标下的 Schwarzschild 光锥。在这些坐标中，指向过去的类时路径可以越过 $r = 2GM$。}

% ---- 书页 224 / p239 ----
这导致
\begin{equation*}
\mathrm{d}s^2 = -\frac{1}{2}\left(1 - \frac{2GM}{r}\right)\left(\mathrm{d}v\,\mathrm{d}u + \mathrm{d}u\,\mathrm{d}v\right) + r^2\,\mathrm{d}\Omega^2,
\tag{5.116}
\end{equation*}
其中 $r$ 通过下式由 $v$ 和 $u$ 隐式定义
\begin{equation*}
\frac{1}{2}(v - u) = r + 2GM\ln\left(\frac{r}{2GM} - 1\right).
\tag{5.117}
\end{equation*}
实际上，我们重新引入了一开始就带着的那种退化；在这些坐标中，$r = 2GM$ 又变得“无限遥远”（位于 $v = -\infty$ 或 $u = +\infty$ 处）。该做的事情是换用一组能把这些点拉回有限坐标值的坐标；一个很好的选择是
\begin{align*}
v' &= e^{v/4GM}\\
u' &= -e^{-u/4GM},
\tag{5.118}
\end{align*}
用我们原来的 $(t, r)$ 坐标系表示，即
\begin{align*}
v' &= \left(\frac{r}{2GM} - 1\right)^{1/2}e^{(r+t)/4GM}\\
u' &= -\left(\frac{r}{2GM} - 1\right)^{1/2}e^{(r-t)/4GM}.
\tag{5.119}
\end{align*}
在 $(v', u', \theta, \phi)$ 坐标系中，Schwarzschild 度规为
\begin{equation*}
\mathrm{d}s^2 = -\frac{16G^3M^3}{r}e^{-r/2GM}\left(\mathrm{d}v'\mathrm{d}u' + \mathrm{d}u'\mathrm{d}v'\right) + r^2\,\mathrm{d}\Omega^2.
\tag{5.120}
\end{equation*}
终于，$r = 2GM$ 的非奇异性变得完全显露出来；在这种形式下，在事件视界处没有任何度规系数表现出任何特殊的行为。

$v'$ 和 $u'$ 都是类光坐标，意思是它们的偏导数 $\partial/\partial v'$ 和 $\partial/\partial u'$ 是类光矢量。这本身并没有任何问题，因为这个坐标系中四个偏导数矢量（两个类光、两个类空）完全可以充当切空间的一组良好的基。尽管如此，在一个坐标类时、其余坐标类空的坐标系中工作还是要让人安心一些。于是我们定义
\begin{equation*}
T = \frac{1}{2}(v' + u') = \left(\frac{r}{2GM} - 1\right)^{1/2}e^{r/4GM}\sinh\left(\frac{t}{4GM}\right)
\tag{5.121}
\end{equation*}
以及
\begin{equation*}
R = \frac{1}{2}(v' - u') = \left(\frac{r}{2GM} - 1\right)^{1/2}e^{r/4GM}\cosh\left(\frac{t}{4GM}\right),
\tag{5.122}
\end{equation*}

% ---- 书页 225 / p240 ----
用它们表示，度规变为
\begin{equation*}
\boxed{\,\mathrm{d}s^2 = \frac{32G^3M^3}{r}e^{-r/2GM}\left(-\mathrm{d}T^2 + \mathrm{d}R^2\right) + r^2\,\mathrm{d}\Omega^2,\,}
\tag{5.123}
\end{equation*}
其中 $r$ 由下式隐式定义
\begin{equation*}
T^2 - R^2 = \left(1 - \frac{r}{2GM}\right)e^{r/2GM}.
\tag{5.124}
\end{equation*}
坐标 $(T, R, \theta, \phi)$ 被称为\textbf{Kruskal 坐标}（Kruskal coordinates），有时也称为 Kruskal--Szekeres 坐标。

Kruskal 坐标具有许多奇妙的性质。与 $(t, r^*)$ 坐标一样，径向类光曲线看起来与平直空间中相同：
\begin{equation*}
T = \pm R + \text{constant}.
\tag{5.125}
\end{equation*}
然而与 $(t, r^*)$ 坐标不同的是，事件视界 $r = 2GM$ 并不无限遥远；实际上它由下式定义
\begin{equation*}
T = \pm R,
\tag{5.126}
\end{equation*}
这与它是一个类光曲面相一致。更一般地，我们可以考虑 $r =$ 常数的曲面。由式 (5.124)，这些曲面满足
\begin{equation*}
T^2 - R^2 = \text{constant}.
\tag{5.127}
\end{equation*}
于是，它们在 $R$-$T$ 平面上呈现为一条条双曲线。此外，$t =$ 常数的曲面由下式给出
\begin{equation*}
\frac{T}{R} = \tanh\left(\frac{t}{4GM}\right),
\tag{5.128}
\end{equation*}
这定义了过原点、斜率为 $\tanh(t/4GM)$ 的直线。注意，当 $t \to \pm\infty$ 时，式 (5.128) 变得与式 (5.126) 相同；因此 $t = \pm\infty$ 与 $r = 2GM$ 代表同一个曲面。

我们的坐标 $(T, R)$ 应当被允许取遍它们在碰不到 $r = 0$ 处真实奇点的前提下所能取的一切值；因此允许的区域是
\begin{align*}
-\infty &\leq R \leq \infty\\
T^2 &< R^2 + 1.
\tag{5.129}
\end{align*}
由式 (5.121) 和 (5.122)，当 $r < 2GM$ 时 $T$ 和 $R$ 似乎会成为虚数，但这是一种错觉；在那个区域里 $(r, t)$ 坐标本来就不好用（具体地说，$|t| > \infty$）。现在我们可以在 $T$-$R$ 平面上（隐去 $\theta$ 和 $\phi$）画出一幅时空图，即所谓的\textbf{Kruskal 图}（Kruskal diagram），如图 5.12 所示。图上的每个点都是一个二维球面。这幅图表示的就是 Schwarzschild 几何的最大延拓，

% ---- 书页 226 / p241 ----
%FIG{5.12}: {Kruskal 图——Kruskal 坐标下的 Schwarzschild 解，其中所有光锥都张开 $\pm 45^{\circ}$。}

这些坐标所覆盖的，正是我们应当看作由这个解所描述的整个流形。

原来的 Schwarzschild 坐标 $(t, r)$ 只在 $r > 2GM$ 时好用，而这只是 Kruskal 图上所描绘的流形的一部分。把图分成四个区域会带来方便，如图 5.13 所示。区域 I 对应于 $r > 2GM$，即我们原来的坐标在其中具有良好定义的那一片。沿着指向未来的类光射线，我们到达区域 II；沿着指向过去的类光射线，则到达区域 III。如果当初探索的是类空测地线，就会被引到区域 IV。把 $(T, R)$ 与 $(t, r)$ 联系起来的定义 (5.121) 和 (5.122) 实际上只在区域 I 中好用；在其他区域中，必须引入适当的负号，以防止坐标变成虚数。

%FIG{5.13}: {Kruskal 图的各个区域。}

% ---- 书页 227 / p242 ----
把 Schwarzschild 几何延拓到所能延拓的极致之后，我们描述了一个不同凡响的时空。当然，区域 II 就是我们所说的黑洞。任何东西一旦从区域 I 进入区域 II，就再也回不来了。实际上，区域 II 中每一条指向未来的路径最终都会撞上 $r = 0$ 处的奇点；一旦进入事件视界，你就在劫难逃。这一点值得强调：你不仅无法逃回区域 I，甚至无法让自己停止沿 $r$ 减小的方向运动，因为这恰恰是类时方向。这一点本来也可以在原来的坐标系中看出来：当 $r < 2GM$ 时，$t$ 变成类空的，而 $r$ 变成类时的。因此，你无法停止向奇点运动，就像你无法停止变老一样。由于固有时沿测地线取极大值，如果你不作挣扎，而是放松身体接近奇点，你就能活得最久。倒不是说你还有多久可以放松，这趟旅程也谈不上令人松弛：当你接近奇点时，潮汐力会变得无穷大。当你落向奇点时，你的脚和头会被彼此拉开，而你的躯干则被挤压到无穷薄。一位天体物理学家落入黑洞后的凄惨下场，在 Misner、Thorne 与 Wheeler (1973) 的第 32.6 节中有细致的描述。注意他们使用了正交归一标架，正如我们在附录 J 中讨论的那样（不过这可不会让这趟旅行变得更愉快）。

区域 III 和区域 IV 可能有些出乎意料。区域 III 不过是区域 II 的时间反演，是时空中这样一个部分：其中的东西可以逃到我们这里，而我们却永远到不了那里。它可以被看作一个\textbf{白洞}（white hole）。在过去存在一个奇点，宇宙看起来正是从那里迸发而出的。区域 III 的边界是过去事件视界，而区域 II 的边界是未来事件视界。至于区域 IV，无论沿时间向前还是向后，都无法从我们的区域 I 到达那里，那边的人也同样无法到我们这里来。它是时空的另一个渐近平直区域，是我们这个区域的镜像。可以认为它通过一个虫洞（wormhole，或称 Einstein--Rosen 桥）与区域 I 相连——这是一种连接两个不同区域的瓶颈状结构。考虑把 Kruskal 图按 $T =$ 常数的类空曲面切片，如图 5.14 所示。现在我们可以把每个切片的图像画出来，为清晰起见把一个角坐标也画上，如图 5.15 那样。按这种切片方式，Schwarzschild 几何描述的是两个渐近平直的区域：它们相互伸向对方，经由虫洞连接一段时间，然后再断开。但虫洞闭合得太快，任何类时观者都来不及从一个区域穿过它进入另一个区域。

尽管 Kruskal 图已经令人赏心悦目，通过构造共形图把 Schwarzschild 解压缩到有限区域之内，往往还要更有用。共形图的思想在附录 H 中讨论；它是分析广义相对论中时空的关键工具，建议你现在就把那里的讨论复习一遍。我们不会完整详尽地给出构造 Schwarzschild 共形图所必需的那些操作，因为它们与 Minkowski 情形如出一辙，只是代数上额外复杂了不少。我们会从类光版本的 Kruskal 坐标出发，在其中度规

% ---- 书页 228 / p243 ----
%FIG{5.14}: {Kruskal 坐标中的类空切片。}

取如下形式
\begin{equation*}
\mathrm{d}s^2 = -\frac{16G^3M^3}{r}e^{-r/2GM}\left(\mathrm{d}v'\mathrm{d}u' + \mathrm{d}u'\mathrm{d}v'\right) + r^2\,\mathrm{d}\Omega^2,
\tag{5.130}
\end{equation*}
其中 $r$ 通过下式隐式定义
\begin{equation*}
v'u' = -\left(\frac{r}{2GM} - 1\right)e^{r/2GM}.
\tag{5.131}
\end{equation*}
然后，基本上照搬平直时空情形中用过的同一个变换，就足以把无限远收进有限的坐标值：
\begin{align*}
v'' &= \arctan\left(\frac{v'}{\sqrt{2GM}}\right)\\
u'' &= \arctan\left(\frac{u'}{\sqrt{2GM}}\right),
\tag{5.132}
\end{align*}

%FIG{5.15}: {图 5.14 中诸类空切片的几何。}
